%Version 3.1 December 2024
% Structured for Overleaf using sn-jnl, following the same conventions as the
% MIMIC-MJX manuscript template. Figures are left as placeholder \includegraphics
% calls pointing at figs/fig1.pdf ... figs/fig6.pdf -- swap in the real PDFs later.
%
%%%%%%%%%%%%%%%%%%%%%%%%%%%%%%%%%%%%%%%%%%%%%%%%%%%%%%%%%%%%%%%%%%%%%%
%% Submit your LaTeX manuscript as one .tex document.                %%
%% All additional figures/files should be attached separately and    %%
%% not embedded in the \TeX\ document itself.                        %%
%%%%%%%%%%%%%%%%%%%%%%%%%%%%%%%%%%%%%%%%%%%%%%%%%%%%%%%%%%%%%%%%%%%%%%

\documentclass[pdflatex, sn-nature]{sn-jnl}
% Center the text block on each page (avoids the twoside binding-offset shift)
% Also increase \footskip so page numbers don't collide with bottom floats/captions.
\geometry{hcentering,bindingoffset=0mm,footskip=150pt}

%%%% Standard Packages
\usepackage{graphicx}%
\usepackage{multirow}%
\usepackage{amsmath,amssymb,amsfonts}%
\usepackage{amsthm}%
\usepackage{mathrsfs}%
\usepackage[title]{appendix}%
\usepackage{xcolor}%
\usepackage{textcomp}%
\usepackage{manyfoot}%
\usepackage{float}%
\usepackage{placeins}%  % \FloatBarrier, keeps floats within their section
\renewcommand{\topfraction}{0.9}%
\renewcommand{\bottomfraction}{0.8}%
\renewcommand{\textfraction}{0.07}%
\renewcommand{\floatpagefraction}{0.7}%
\setcounter{topnumber}{2}%
\setcounter{bottomnumber}{2}%
\setcounter{totalnumber}{4}%
\usepackage{bm}%
\usepackage{booktabs}%
\usepackage{algorithm}%
\usepackage{algorithmicx}%
\usepackage{algpseudocode}%
\usepackage{listings}%
\usepackage{caption}
\usepackage{tabularx}
\usepackage{makecell}
\usepackage{array}

\captionsetup{font=footnotesize}

\theoremstyle{thmstyleone}%
\newtheorem{theorem}{Theorem}%
\newtheorem{proposition}[theorem]{Proposition}%

\theoremstyle{thmstyletwo}%
\newtheorem{example}{Example}%
\newtheorem{remark}{Remark}%

\theoremstyle{thmstylethree}%
\newtheorem{definition}{Definition}%

\raggedbottom

\begin{document}

\title[Article Title]{LUC3D: Browser-based 3D annotation and proofreading GUI for pose estimation and multi-animal tracking}

%%=============================================================%%
%% GivenName -> \fnm{}   Particle -> \spfx{}   FamilyName -> \sur{}
%%=============================================================%%

\author[1]{\fnm{Eric J.} \sur{Leonardis}}
\author[2]{\fnm{Stefan N.} \sur{Oline}}
\author[1]{\fnm{Joshua} \sur{Park}}
\author[2]{\fnm{Bartul} \sur{Mimica}}
\author[1]{\fnm{Liezl} \sur{Maree}}
\author[2,1]{\fnm{Shayan} \sur{Afshar}}
\author[1]{\fnm{Isaac} \sur{Tang}}
\author[1]{\fnm{Kay M.} \sur{Tye}}
\author[2]{\fnm{Mala} \sur{Murthy}}
\author[2]{\fnm{Annegret L.} \sur{Falkner}}
\author*[1]{\fnm{Talmo D.} \sur{Pereira}}\email{talmo@salk.edu}

\affil[1]{\orgname{Salk Institute for Biological Studies}, \orgaddress{\city{La Jolla}, \state{CA}, \country{USA}}}
\affil[2]{\orgname{Princeton Neuroscience Institute, Princeton University}, \orgaddress{\city{Princeton}, \state{NJ}, \country{USA}}}

% max 150 words, unreferenced
\abstract{Animal behavior unfolds in three dimensions over time, yet accessible scientific tools for analyzing it in this native form remain scarce. Uncertainty from 2D pose methods can be refined with multi-camera approaches to allow for more accurate 3D estimates. This paper presents LUC3D (Label Unification and Correspondence in 3D Annotation GUI), a browser-based 3D GUI for reducing the amount of manual labor required for labeling, proofreading, and multi-camera multi-animal tracking while requiring no installation to maximize accessibility. In addition, we present multiple comprehensive datasets of proofread multi-camera 3D triangulated pose data of multi-animal mouse social behavior. Also, a protocol for generating 3D consistent labels is presented and a cross-view ReID tracking algorithm for proofreading multi-animal multi-camera pose data is made available. This tool will allow for the widespread proliferation of 3D pose estimation techniques for studying social behavior while reducing the cost of entry for neuroscience labs across the world.}

\maketitle

\section{Introduction}\label{sec1}

Advances in deep learning for 2D pose estimation in humans has led to the development of accessible and popular methods for quantifying animal behavior \citep{Mathis2018,Pereira2020,Pereira2022,Biderman2024}. The human pose estimation literature has demonstrated that the solution to these myriad problems caused by the limitations of 2D is to use multiple camera views to enable 3D pose estimation \citep{Iskakov2019}. This methodology has been applied to 3D pose estimation of animal behavior \citep{Dunn2021,Karashchuk2021}. Specialized tools have also been created for tracking the 3D movement of the face in rodents \citep{Daruwalla2026}. These single-subject 3D methods have since been extended to enable 3D pose estimation of multiple individuals simultaneously \citep{Huang2020,Chen2020}, a methodology that has similarly been extended to multi-animal applications \citep{Marshall2021,Marshall2022}. While other 3D annotation GUIs exist, there can often involve a difficult installation process which is non-trivial \citep{Hueser2023}. Recently LightningPose3D has been developed for multi-view 3D pose estimation and labeling using browser-based applications that require an additional installation \citep{Aharon2026}. 

Here we present LUC3D, a browser-based GUI that brings together SLEAP's to allow for 3D consistent cross-view labeling, tracking, and proofreading without any installation required whatsoever to maximize accessibility. The LUC3D GUI allows for the user to leverage 3D triangulation and reprojection to streamline the labeling process for multiple camera views with careful consideration of user experience. Reprojection-aided labeling provides the user with the opportunity to speed up the labeling process by using reprojections from two anchor views into the remaining views to save the user time from clicking on each node separately. After reprojection the user can either accept or modify the reprojections to best reflect the pose in the remaining views. This tool ships with a cross-view tracker that uses 3D information to allow for the identification of the same animal across different cameras \citep{Maree2024}. We also provide a browser-based tool for camera calibration to ensure that users have all the necessary tools for 3D pose estimation with no installation necessary. 

In addition to new software tools, this paper presents two comprehensive datasets: the Mouse-Dyad-10M two-mouse data repository and the SLAP-2M data repository which has up to 4 mice and multiple levels of environmental enrichment. Many popular 2D datasets lack environmental enrichment because they introduce occlusions, by adding levels of enrichment this dataset highlights the need for 3D pose estimation. This makes the SLAP-2M dataset an ideal candidate for benchmarking new models for behavioral segmentation 3D pose estimation because it has different levels of difficulty determined by number of enrichment objects and number of animals.

Several practical issues make 3D multi-animal pose estimation harder to adopt than single-camera methods. Building a multi-camera training set can be costly, since single-camera pose estimation methods require many manually labeled frames per view before a model performs reliably. Reprojection-aided labeling lets a frame labeled in a few anchor views generate accurate labels in all remaining views, sharply reducing the number of manual labels needed to generate a consistent training set. Proofreading requires requires resolving identity separately within each view before identity can be linked across views. The LUC3D cross-view tracker resolves identity across all cameras at once, so proofreading a session requires reviewing it once rather than once per camera. Adding more cameras also gives the tracker more redundant evidence for improving identity estimates. This paper provides a browser-based 3D annotation, proofreading, tracking, and calibration interface that makes 3D multi-animal pose estimation more usable and accessible. We also provide two massive datasets that serve as challenging benchmarks for future 3D cross-view tracking methods.

\section{Results}\label{sec2}

\subsection{A browser-based tool that speeds up annotation and resolves cross-view identity}\label{subsec2-1}

Multi-camera pose recording poses two main problems: pose tracking within each camera view and linking identities across them. Figure~\ref{fig1} A shows two multi-view recording setups: the SLAP-2M cage with eight cameras and up to 4 animals, and the Mouse-Dyad-10M arena with five cameras and 2 animals, rendered from a single session's own calibration and tracked 3D poses. The first problem, finding the keypoints in each view, is solved by existing detectors. The second, deciding which detection in one view is the same animal as which detection in another, which currently has no standard tool, and until it is solved there is cannot be 3D reconstruction and nothing to proofread. Usually, the cross-view ID problem is solved by manual assignment or assuming that the initial tracks have no dropped frames whatsoever \citep{Klibaite2025, Waldmann2024}. The three stages it contributes are cross-view re-identification, triangulation, and proofreading of the resulting 3D building upon the 2D pose innovations from SLEAP \citep{Pereira2022}.

LUC3D addresses the problems presented by multi-camera 3D annotation and tracking in a web browser, with no installation and no build step (Figure~\ref{fig1}B). It contributes three stages: cross-view re-identification, triangulation, and proofreading of the resulting 3D. The per-view detector comes from elsewhere; the application consumes SLEAP predictions. Figure~\ref{fig1}C shows the operation on a single frame of an eight-camera recording of three mice. The detector supplies 25 detections across the eight views, carrying 21 distinct track names, and those names are per camera and arbitrary: the same animal is track 89 in one view and track 83 in another. After re-identification the frame holds three identities, one per animal in every view, with 24 of the 25 detections assigned. The one left over is a duplicate detection of an already-matched animal in a view the panel does not show, which one-to-one assignment correctly refuses. Figure~\ref{fig1}D triangulates that frame and fills all 45 of its 3D keypoints, placing the reconstruction at one camera's own pose and field of view so that it can be compared directly against that camera's video. Figure~\ref{fig1}E places the tool among seven existing packages. Reprojection-aided labeling predates this work: Label3D and JARVIS both provide it \citep{Hueser2023, Aldarondo2019}. No other tool in the table combines annotation-time cross-view identity, multi-animal support, a 3D proofreading viewport, and running in a browser without installation.

\begin{figure}[p]
  \centering
  % Keep this figure+caption away from the footer/page number region.
  \includegraphics[width=1.0\textwidth,height=1.0\textheight,keepaspectratio]{figs/fig1.pdf}
  \caption{\textbf{LUC3D annotates and proofreads multi-camera 3D pose in a browser with no installation, and assigns cross-view identity} \textbf{A}, The SLAP-2M rig, left: 8 calibrated cameras and the cage with the tracked 3D poses of four animals. The Mouse-Dyad-10M rig, right: the 5 calibrated cameras and enclosure with tracked 3D poses of the two animals. \textbf{B}, Pipeline, from videos to export. Video acquisition, 2D pose estimation (SLEAP or any other predictor as .slp), cameras are calibrated with a browser-based tool, cross-view re-identification, one per animal; triangulation collapses the tiles into a single 3D volume; the 3D is proofread against per-view reprojections; and the project exports as .slp 2.8 or HDF5. The bracket marks the stages contributed by this paper. \textbf{C}, One frame of an 8-camera, 3-mouse recording, 15-node skeleton. Left, the per-camera tracks supplied to the app; right, the same two views after cross-view re-identification (cam 0 mid and cam 7 sideR, 2 of 8 views). Labels are the per-camera track name and the resolved identity. Track colors are arbitrary within each camera and carry no meaning across views; identity colors are shared by every view. Across all 8 views this frame holds 25 detections carrying 21 distinct track names; 24 of the detections are assigned, resolving to 3 identities, one per animal in every view, and 1 detection is left unassigned. \textbf{D}, The same frame triangulated. Left, cam 0 mid with identity overlays; middle, the 3D viewport placed at that camera's own pose and field of view; right, the rig, with all 8 calibrated cameras and the 3 reconstructed animals. 45 of 45 3D keypoints filled. \textbf{E}, Capability comparison against 7 multi-animal or multi-camera pose tools, read from each tool's published documentation. A dash means the capability is not documented or not applicable.}
  \label{fig1}
\end{figure}

\subsection{Buy two get N free: reprojection speeds up the labeling process}\label{subsec2-2}

Traditionally, manual multi-camera annotation scales with the number of cameras. At the initial stages of labeling, the labeler places every keypoint in every view, so the cost of a frame grows with the rig. If the 3D point can be solved from two views, the remaining views can be filled by reprojection, and the labeler only has to correct those reprojections that lack good correspondence to the animals in the video. When using reprojection-aided labeling, the labeler can save time by accepting good reprojections with one click and then only having to modify the fraction of reprojections that do not correspond to the animals in the video. 

Figure~\ref{fig2}B examines the required number of placements per animal per frame for manual labeling in comparison with reprojection-aided labeling. On the measured five-camera rig, accepting reprojections at 10 px reduces manual placements from 75 per animal to 32, a 2.3 reduction in manual labeling. While modifying those reprojections to reduce the error to approximately 5 px would result in 48 placements per animal, which is a 1.6-fold reduction in manual labeling. The saving grows with rig size because the labeled two views stay fixed while the reprojected remainder does not. This shows an estimate of the amount of labeling effort reduced by using reprojections in the labeling process. This assumes that each node would be manipulated separately, but the GUI provides methods for translating the whole instance which can also reduce labeling time. 

The relationship between the numbers of cameras used in triangulation and the 3d error demonstrates that multi-view geometry reduces the uncertainty of pose estimates. A 3D point solved from two views and projected into a view that was not labeled lands a median 4.27 mm and the five-view solve 1.19~mm (Figure~\ref{fig2}C). The reference is proofread 3D reconstruction aligned to the calibration frame. 

The cost of solving from only two views depends on where those two cameras are placed relative to each other Figure~\ref{fig2}D. Error therefore falls as $k$ over the sine of the angle between the two cameras relative to the animal (dashed curve).  For example, the narrowest angle pair on this rig at 13 degrees gives 12.6~mm of error while the widest at 31 degrees gives 2.7~mm, against a floor of 1.2~mm when all five views contribute. So this means that strategic camera placement when setting up the rig and strategic choice of the two anchor views for reprojection can also reduce labor cost. Increasing the number of contributing views leads to more accurate 3D estimates compared with ground truth, scored relative to a camera the solve never saw, error falls from 4.32~px with two cameras to 3.34~px with four views (Figure~\ref{fig2}E (Top)). Figure~\ref{fig2}E (Bottom) shows that dropping an unreliable view can reduce reprojection error.

\begin{figure}[!htb]
  \centering
  \includegraphics[width=1\textwidth]{figs/fig2.pdf}
  \caption{\textbf{Reprojection-aided labeling protocol.} \textbf{A}, The protocol in the luc3d GUI: two anchor views labeled (cam 1 topB, cam 6 sideL), the 3D solved from those two alone, the reprojection drawn into the unlabeled views (cam 0 mid, cam 2 topC), and one view magnified beside that frame's per-view error split. \textbf{B}, Manual placements per animal per frame against rig size $C$: every label by hand ($C \times N$, $N = 15$ nodes) against the placements still needed after reprojection at $\tau = 10$~px (solid) and 5~px (dashed). Curves are the model; the filled marker is the measured rig, $C = 5$, where 32 of 75 are placed by hand, 2.3-fold fewer. \textbf{C}, 3D error against the cameras in the solve. Boxes, the 50 session medians (median, IQR, whiskers 1.5$\times$IQR), every session also a dot: 4.74, 2.89, 1.91 and 1.19~mm at 2 to 5 cameras. \textbf{D}, Median two-anchor 3D error against the angle the pair subtends at the animal, one marker per camera pair. Dashed curve, $k/\sin\theta$ with $k = 1.52$~mm; band, plus or minus 25\%; dotted line, the all-five-view floor, 1.2~mm. 12.6~mm at 13.5 degrees falls to 2.7~mm at 31.5. Mouse-Dyad-10M: B and D, all 56 sessions, 286,200,174 keypoints; C, the 50 proofread sessions. A is one frame of an 8-camera Hard-Fight session. \textbf{A}, Triangulation accuracy is set by how many views contribute and by whether a badly fitting view is dropped, not by the solver. The two solvers the app ships: linear DLT, closed form, the default; and non-linear triangulation in native pixels with the cameras fixed (app menu: ``Bundle Adjustment''). Padlocks mark fixed cameras. \textbf{e}, Error in a camera outside the solve, every subset at each count: DLT 4.32, 3.66 and 3.34~px and refinement 4.42, 3.53 and 3.15~px at 2, 3 and 4 cameras. Bars, the distribution-free 95\% CI of the across-session median; paired, refinement minus DLT is $+0.111$, $-0.069$ and $-0.098$~px. Reprojection error in the kept views, all views in the solve against the same solve with the worst view dropped; one line per session, bars the 95\% CI of the mean. 2.06 to 1.71~px, a paired $-0.345$~px (95\% CI $-0.359$ to $-0.332$), lower in 50 of 50 sessions. \textbf{F}, Median reprojection error per session in the cameras the solve used, four solvers paired by class. Boxes, the session distribution (median, IQR, whiskers 1.5$\times$IQR): Anipose linear 2.26, LUC3D DLT 2.35, Anipose optim 2.26 and LUC3D refined 2.15~px, lowest in 50 of 50 sessions. \textbf{G}, Solve time per keypoint, undistortion excluded from all four: 6.3 against 29.0~\textmu s on the linear pair, 44.0 against 228.8~\textmu s on the non-linear. $n = 50$ Mouse-Dyad-10M sessions at every 15th frame, 17,013,412 keypoints, the same keypoints in every column. Anipose is aniposelib 0.7.2, ransac and smoothing off.}
  \label{fig2}
\end{figure}

\subsection{A greedy per-camera assignment outperforms exhaustive approaches and enables fast tracking}\label{subsec2-3}

Maree et al. (2024) \citep{Maree2024} uses a cross-view association algorithm that enumerates every possible grouping of detections into identities, triangulates and reprojects each, and keeps the grouping with the lowest reprojection error. That is $A!$ per view and therefore $(A!)^C$ per frame: 32 hypotheses for two animals in five cameras and 190 million for four animals in six (Figure~\ref{fig3}B). Following the approach of Chen et al. (2020) \citep{Chen2020}, LUC3D instead solves one assignment problem per camera and commits it before moving to the next, at a cost of $C \cdot A^3$ (Figure~\ref{fig3}A).

We investigated whether the cheap and greedy procedure gives up accuracy. Figure~\ref{fig3}E answers it against proofread ground truth rather than against the other method, over 92 sessions: the 50 Mouse-Dyad-10M sessions at two animals in five cameras, plus the 42 SLAP-2M sessions that carry both pool detections and proofread ground truth. 

Figure~\ref{fig3}F shows the time it takes to run the two procedures on identical detections per frame. The gap is a factor of about a million in the four-animal, six-camera configuration, which exhaustive enumeration could not run at all. That point is drawn as a bound rather than a measurement, and the bound is generous to the published method twice over: it removes a relabeling symmetry that method does not exploit, and it prices the remainder at the cheapest rate measured anywhere in the sweep. Even so, one second of 50~fps video would cost more than a day.

Two limits belong with this result. Exhaustive enumeration is undefined unless every camera holds exactly as many detections as there are animals, so it was run only on the 47.4 per cent of frames that satisfy that, 4,591,864 of 9,678,503, and the excluded frames are the occluded ones, which is to say exactly the frames association finds hardest. And 94.2 per cent of the computed frames are the easiest configuration, two animals in five cameras. Agreement between the two methods falls monotonically as animals are added, from 99.996 per cent at two animals in five cameras to 99.038 per cent at four animals in three, a trend the earlier four-session version of this benchmark could not have shown. The equivalence therefore holds where the published method can run, and is not a claim about the frames where it cannot.

Figures~\ref{fig3}G and H ask what the three-dimensional term in the association cost contributes, by sweeping its weight against the two-dimensional term's over all 50 Mouse-Dyad-10M sessions, in the same configuration as Figure~\ref{fig3}E. The two weights matter only through their ratio, a fact verified on an earlier eight-session grid where all 24 weight combinations collapsed onto it, so the sweep samples the ratio directly. The term is necessary. The association does not survive without it: with the term switched off entirely, identity switches run at 632 per 100,000 camera-frames and cross-view IDF1 falls to 0.599. At the app default ratio of six the arm holds 413 switches, 0.92 per 100,000, at IDF1 0.8613, and doubling the ratio to twelve buys 42 fewer switches and no IDF1 (0.8614 against 0.8613). Rates are over the 45,021,960 camera-frames of the 50 sessions.

For further characterization of the LUC3D cross-view tracker performance in comparison with other trackers, like SLEAP and ByteTrack, and for additional parameter searches which reduce ID switches (see Supplementary Figure~\ref{fig6} and Supplementary Sections~\ref{subsec2-7} and~\ref{subsec2-9}).

\subsection{Triangulation can be performed quickly without giving up quality}\label{subsec2-4}

We have demonstrated that there is little difference in the quality of the solvers when comparing Anipose with LUC3D. However LUC3D is able to perform triangulations faster which supports real-time usability in the browser-based GUI. Across four solvers, two from LUC3D and two from aniposelib, the range of per-session medians is betweem 2.15 to 2.35~px (Figure~\ref{fig2}F). Neither library's non-linear refinement shows a large improvement for this dataset. A speed benchmark test showed that LUC3D's solvers are 4.6 and 5.2 times faster per keypoint than the corresponding aniposelib paths (Figure~\ref{fig2}G) due to a smaller number of iterations.

\begin{figure}[!htb]
  \centering
  \includegraphics[width=1\textwidth]{figs/fig3.pdf}
  \caption{\textbf{Cross-view tracking algorithm.} Greedy per-camera assignment groups as well as exhaustive enumeration at lower cost. \textbf{A}, The two strategies: exhaustive enumeration of every grouping \citep{Maree2024}, $(A!)^C$ per frame, against LUC3D's one Hungarian assignment per camera, each committed before the next, at $O(C \cdot A^3)$. \textbf{B}, Hypotheses per frame for exhaustive enumeration, one curve per rig size; dotted line, the harness cap of $10^6$. \textbf{C}, An illustration of ID grouping hypotheses between two camera views for 3 animals, incorrect groupings shown by the gray lines and correct groupings shown in color forming the 3D triangulation..  \textbf{D}, An identity switch, staged from one real HardFight frame in the same two-view vocabulary as C. The top camera's two identities are exchanged (the white curved arrows); the side camera is correct. Solid legs tie each correct 2D detection to its 3D instance on the floor; the dotted legs run a colour gradient between the two identities, so a leg that starts orange and ends blue IS the statement that one animal is carrying two identities. \textbf{E}, Grouping accuracy: frames misgrouped against proofread ground truth per 100,000 clean frames, pooled across the animals-x-cameras configurations; boxes, the per-session distribution on a symlog axis, every session a dot, the white line the median. LUC3D's greedy assignment misgroups 926 frames against exhaustive's 1,309 over 4,572,172 clean frames (20.3 against 28.6 per 100,000); at 4 animals in 3 cameras the per-session medians are 0 against 880 per 100,000. \textbf{F}, Time per frame on identical detections: exhaustive 11.5 ~ms at 2 animals in 5 cameras to 1,980.9 s at the 4-in-6 lower bound, LUC3D 1.1 to 2.4 ~ms. \textbf{G}, Parameter sweep of the 3D term on the ratio $r = corr3d/corr2d$: switches per 100,000 camera-frames for exhaustive vs. greedy (log). \textbf{H}, Parameter sweep of the 3D term on the ratio $r = corr3d/corr2d$: cross-view IDF1 for exhaustive vs. greedy. n = 92 sessions (50 Mouse-Dyad-10M, 42 SLAP-2M); 4,591,864 of 9,678,503 frames were clean and all were computed. C is one frame of a 3-animal SLAP-2M session, D one frame of the HardFight recording; G and H are Mouse-Dyad-10M.}
  \label{fig3}
\end{figure}

\subsection{Cross-view identity tracking enables the analysis of 3D social behaviors}\label{subsec2-5}

The preceding figures demonstrate the usefulness of 3D multi-animal pose estimation data for understanding social behavior. Figure~\ref{fig4} measures a social rearing behavior in the Mouse-Dyad-10M corpus. This social rearing behavior has leader and follower dynamics that may be related to courtship behavior between males and females. 

Rearing is defined increasing head height above the floor on the two hind legs, often with the forelimbs supporting the animals weight on a wall. We were particularly interested in the social case when rearing shows up as a synchronized pair interaction. In the example display of Figure~\ref{fig4}A every camera sits between 58 and 76 degrees above the animals, and in every one of the five views the two noses sit 2.7 to 6.0~px apart while the tail bases sit 93 to 106~px apart. No single view recovers the vertical axis due to the lack of a z axis measurement and no view turns can turn pixel gaps into a metric distance between the animals.

Rearing in these pairs is coupled and the coupling requires proximity. When averaging the likelihood of rearing given the onset of the other animal demonstrates an asymmetric relationships (Figure~\ref{fig4}G). Split by which animal's onset anchors the window, the male's rearing near the female's onset is below the male's own chance rate at the instant the female starts (0.6 times chance) and reaches only 1.7 times about 0.8~s later. While the female's rearing near the male's onset is already 4.7 times chance at the instant he starts, peaking at 5.3 times a third of a second later. For the male's near-rear onsets, the female is already mid-rear 50.2 per cent of the time. So the male's onset is usually joining a rear the female has already started. For the female's near-rear onsets, the male is already mid-rear only 6.2 per cent of the time, so the female's onset is usually a fresh start. Neither half of the split is a shared external drive, since both require proximity in the same way the pooled measurement does. The same measurement on the two-animal sessions of SLAP-2M gives 1.08 within two body lengths and 0.97 beyond, that is no coupling at all, and no systematic sex convention exists there to split by. Arena size may also be a likely reason for the SLAP-2M null.

That proximity coupling defines an event, hereafter referred to as a social rearing display. During these displays, both animals reared within two body lengths and held for at least a quarter of a second. There are 539 such displays in 37 of the 56 sessions, with a median duration of 0.71~s. In the half second immediately before a display begins the two sexes are already moving very differently (Figure~\ref{fig4}D): the male is still traveling, at a median 0.38 body lengths per second, while the female is already largely stationary at a median 0.23 (paired Wilcoxon $P = 6.5 \times 10^{-22}$, $n = 538$ displays with both speeds resolved). The difference is not only in how fast each animal moves but in what that motion is doing. Decomposing each animal's own body orientation towards its partner and scaling by that animal's own speed (Figure~\ref{fig4}E), the male's score in the 1.5~s before onset is strongly negative (median $-0.71$: the male's body axis is oriented at the female's and he is moving), while the female's is close to zero (median $-0.06$; paired Wilcoxon $P = 3.3 \times 10^{-73}$). The male is oriented at the female's body and closing the distance; the female is neither oriented at him nor approaching him.

The sequential ordering of the display is not symmetric and the leader-follower dynamic is sexually dimorphic. The female is up a median 0.39~s before the male joins, in the 432 of 539 displays the female leads (Figure~\ref{fig4}C); at this recording rate that is 59 frames, and the shorter lags in the distribution, down at the 0.17~s quartile, are where the high frame rate earns its keep. The female starts most of the displays. Pooled over all 539 displays the female starts 432 of them, 80.1 per cent and the female is the outright leader starting more displays than the male in 36 of the 37 sessions with at least one display. The one exception is an exact tie, and no session has the male as the outright leader. Because the female's share and the male's share are complementary within a session (they sum to one), Figure~\ref{fig4}F shows that over the 23 sessions with at least six displays, the female's median share is 0.86 against the male's 0.14 (paired Wilcoxon $P = 2.7 \times 10^{-5}$).

The 56 Mouse-Dyad-10M sessions are repeated recordings of 9 mice in 18 pairings rather than 56 independent samples, so the analysis was repeated with the pair as the unit of replication. Every one of the 17 distinct pairings that produced at least one display has the female as its leading member. Restricting to the 14 pairs with at least five displays, covering 536 displays, the female starts a median 0.81 of that pair's displays and leads in all 14 of 14 (sign test $P = 1.2 \times 10^{-4}$), and pooling over pairs gives 429 of 536 displays, 80.0 per cent (the figure the session-level pooling gives). The asymmetry is a property of the pair and of the female in it rather than of any one recording.

Three controls test the obvious alternatives to a genuine sex difference. The female is not simply the animal that rears more, since a session's female-initiation share is uncorrelated with the female's own share of rearing time within the pair ($r = -0.012$, $P = 0.96$ over 24 sessions). It is not an artifact of the per-animal height threshold, since an absolute 60~mm threshold shared by both animals returns the female as leader in all 24 sessions, and a single shared per-pair threshold in 22 of 24. Nor is it explained by body size in the direction that would trivially produce it: the male is the structurally longer animal (median body length 90.1~mm against 84.2~mm, longer in 29 of 37 sessions, paired Wilcoxon $P < 0.0001$), so the smaller-bodied sex is the one leading, and the sign test on the 36 sessions with a unique leader agrees (female shorter in 28 of 36, $P = 0.0012$). Within the same 24-session family as the tests above, the three sessions where the female happens to be the longer animal still hand the female 75 of 97 displays, 77 per cent (binomial $P = 6 \times 10^{-8}$), and over all sessions where the female is the longer animal the count is 81 of 103. Thus these findings demonstrate a sex-based asymmetry for leader-follower dynamics in social rearing that is stable within a session and survives alternative controls. This behavior follows courtship signatures like a female-paced engagement with male pursuit, but cannot be conclusive without further verification of ultrasonic vocalizations, estrous staging or mating assay to accompany these results \citep{Dias2025}.  


\begin{figure}[!htb]
  \centering
  \includegraphics[width=1\textwidth]{figs/fig4.pdf}
  \caption{\textbf{Behavioural analysis of social rearing.} Two mice rear together face to face; the female starts 80\% of the displays and the male joins in. \textbf{A}, One display: the metric 3D reconstruction (230 x 230 x 140 mm; the wall nearest the viewer is drawn as edges only so the interior is seen through clear air) beside the same instant projected into each of the five cameras. Every camera sits 58 to 76 degrees above the animals, so the height that defines the event exists only after triangulation. \textbf{B}, Nose height by sex and the nose gap, around display onset: across-session median of per-session medians, band the across-session p25 to p75. Both animals rise and peak together, the noses converge to 0.12 body lengths at onset, and the female reaches higher (hers is the higher peak on 80.\% of displays -- the curves are keyed by SEX, not by within-display rank). \textbf{C}, How long the female is up before the male joins, over her 432 female-led displays: median 0.39 s (p25 to p75, 0.17 to 0.90 s); hatched bar, the 9\% beyond 2 s. The male-led displays are excluded rather than folded in (their own lag is similar, median 0.32 s). \textbf{D}, Speed in the 0.5 s before onset, over each animal's own session baseline: the male is travelling (median 0.382 body lengths/s) and the female is not (0.231); boxes, median, IQR and 1.5x-IQR whiskers; paired Wilcoxon over the 538 displays with both defined, P = 6.5e-22. \textbf{E}, Facing-pursuit in the same pre-onset window -- each animal's facing axis dotted with the direction away from its partner, scaled by its own relative speed, so negative means moving TOWARD the partner along its own heading: male median -0.708 against female -0.058 (paired Wilcoxon, n = 538, P = 3.3e-73). The male is pursuing; the female is neither approaching nor fleeing. \textbf{F}, Share of displays started, by sex, one value per session: session medians 0.857 (female) against 0.143 (male) over the 23 sessions with six or more displays (paired Wilcoxon P = 2.7e-05); the female leads outright in 36 of 37 sessions with any display (1 tie). Pooled over all 539 displays the female starts 432, 80.1\%; dashed line, the fair coin at 0.5. \textbf{G}, Rear-onset coupling, both directions: the probability the OTHER animal is rearing at each lag around a rear onset, over that animal's own chance rate, for onsets with the pair within 2 body lengths; yellow, the same onsets further apart; grey, a circular-shift null. Around a FEMALE onset (left) the males sits below their chance rate at lag 0 (0.6x, he is not yet up) and climbs to 1.7x about 0.8 s later; around a MALE onset (right) the females are already 4.7x above theirs at lag 0, peaking at 5.3x a third of a second later, the females are already mid-rear at 50.2\% of his near onsets, against 6.2\% for the reverse. Mouse-Dyad-10M: 539 displays in 37 of 56 sessions (2 mice per session, one male and one female, 5 cameras, 150 fps); D and E use the 538 displays with both animals tracked in the pre-onset window; G uses 9,354 rear onsets over all 56 sessions, 2,915 of them within 2 body lengths. A display is both animals reared, neck above 0.75 body length, tail bases within 2 body lengths, held at least 0.25 s.}
  \label{fig4}
\end{figure}

\subsection{Multi-view recovers missing points without sacrificing overall pose quality}\label{subsec2-6}

The SLAP-2M dataset was purposely collected to provide multiple challenges for 2D and 3D pose estimation systems. Each of the 74 SLAP-2M sessions vary by number of animals and number of environmental enrichment objects (see Figure~\ref{fig5}A for a visualization) which increase difficulty due to occlusions. Figure~\ref{fig5}B characterizes how LUC3D tracker's IDF1 score degrades with increased difficulty levels ranging from highly accurate ID tracking with two animals and no enrichment objects (median IDF1 = .989) to multiple animals and multiple enrichment objects (median IDF1 = .654). Figure~\ref{fig5}C shows how increasing the number of camera views can recover missing keypoints. From the easiest stratum to the hardest, measured on every frame of all 74 sessions, the per-view miss rate rises 10.8-fold, from 5.3 to 57.7 per cent. Meanwhile, the mean error of the detections that do fire rises only 1.30-fold, from 3.65 to 4.74~px Figure~\ref{fig5}E (middle). Figure~\ref{fig5}D demonstrates how the animal count affects missing keypoints across difficulty ratings. Since the enrichment objects are transparent and the animals are not, more animals does lead to an increase in keypoints missing within difficulty rating levels. 

\begin{figure}[!htb]
  \centering
  \includegraphics[width=1\textwidth]{figs/fig5.pdf}
  \caption{\textbf{Multi-animal mouse social behavior datasets.} Difficulty removes keypoints rather than degrading them. \textbf{A}, The two axes that set a session's difficulty, rendered from the data. Each cage is one SLAP-2M session recorded at that combination of animal count and environmental enrichment, drawn from its own cage corners and tracked 3D poses. Only the conditions the corpus contains are shown, so the empty cells are absent from the data rather than omitted from the figure. The green enrichment objects are illustrative props at the rated level, since no object positions were recorded. The wall nearest the viewer is drawn as edges only, so the interior is seen through clear air. The inset, marked by the dashed box and leaders, shows the same instant of the 4-animal session as seen by its six proofread cameras, each view cropped to the app's bounding box, with scale bars of 50 mm. An animal carries the same color in the inset as in the cage the inset expands. \textbf{B}, Cross-view IDF1 by difficulty over the 42 multi-animal sessions; boxes, stratum median and IQR, whiskers 1.5x IQR, one dot per session: medians 0.989, 0.917, 0.969, 0.675, 0.877 and 0.654 at ratings 2 to 7. \textbf{C}, Keypoints missing per view after reprojection recovery, by difficulty and rig size: at rating 7, 57.1\% at two cameras falls to 37.7\% at six. \textbf{D}, The same miss rate within each animal count (n = 32, 35, 4, 3 sessions for 1 to 4 animals); bars, +-1 s.d. between sessions; hollow markers, single-session cells. \textbf{E}, Raw per-camera detection quality against difficulty: keypoints missing (10.81-fold), error when present (mean +- s.d. between sessions, with p95; 1.30-fold), and the fraction beyond the app's 20 px tolerance (5.72-fold). \textbf{F}, Per stratum: sessions, keypoints, animal counts, bedding, error mean, p95, p99 and the fraction beyond 20 px. SLAP-2M: C, D, E, F on all 74 sessions; B on the 42 multi-animal sessions.}
  \label{fig5}
\end{figure}

\FloatBarrier

\section{Discussion}\label{sec3}

The objective of LUC3D is to provide an easily accessible interactive GUI for resolving 3D multi-animal annotation and proofreading without any installation required. This tool can speed up the labeling and proofreading process significantly by leveraging reprojection-aided labeling and cross-view identity tracking which would otherwise require significant manual labor. We also provide a valuable resource to the community by sharing two large datasets of multi-animal multi-view 3D pose estimation for examining social behavior: SLAP-2M and Mouse-Dyad-10M. 

We have presented a lightweight and fast method for triangulation using linear and nonlinear methods that allow for significant speedups for the 3D labeling and proofreading process. In addition to the major speed-up also, LUC3D's triangulation performance is comparable to other standard 3D pose estimation methods. Also by providing a browser-based tool for camera calibration we ensure that users have the necessary tools to perform multi-view synthesis without any installation necessary. 

Previous work suggested that exhaustive hypothesis testing would be required to resolve the cross-view identity tracking problem but this becomes untenable when dealing with multiple animals and many camera views \citep{Maree2024}. The results show that greedy Hungarian matching not only can perform this matching efficiently but also can perform better than the exhaustive strategy. This enables real-time interactive cross-view tracking that can be used in a browser-based GUI environment. 

LUC3D utilizes labeling redundancy to reduce labor and minimize manual annotation and provides a method for accepting or modifying reprojections from a small number anchor camera views to the larger number of remaining views. Reprojections can themselves serve as labels, a keypoint labeled across a few views produces a more accurate label in each of those views than labeling any one view alone. This works particularly for occluded nodes, which LUC3D can recover by utilizing reprojection. Reprojection-aided labeling allows the generation of training sets from scratch using far fewer manual labels than a traditional single-view set up.

For multi-animal datasets, correcting 2D identity switches during proofreading is one of the largest remaining blocks of manual labor. Without LUC3D it scales linearly with the number of cameras (a two-camera rig doubles the proofreading time and a four-camera rig quadruples it). Identity must also be associated across views, so that the same animal is recognized as the same track in every camera, which can take additional effort. The LUC3D cross-view tracker instead produces an automated proofread output for both cross-view identity assignment and temporal identity tracking. As a result, the user needs only to review the session once. As a result, proofreading labor is reduced. Additionally, using more cameras give the LUC3D cross-view tracker more redundant evidence so it makes fewer errors that require manual correction. Also the GUI provides visual aids and tools which make the errors easy to catch and correct. 

We have also presented an example of how LUC3D can be used to recover 3D social behavior in a multi-animal context for downstream behavioral analysis. This analysis reveals that rearing in the environment may rely on social contingencies that may be related to courtship or more generally joint exploration behavior \citep{Dias2025, Leonardis2022}. The social rearing analysis also surfaced a sex-specific asymmetry that would not be visible without cross-view identity and metric distance. The female consistently leads the display, starting the large majority of them across nearly every session and every pairing. The male's approach in the moments beforehand looks like directed pursuit, the male stays oriented toward and closing on the female while the female has already stopped moving. The female holding her position and male closing the distance before the pair rears together, resembles a courtship interaction, though we are careful not to overstate this. Confirming the courtship interpretation would require further investigation with independent measures such as ultrasonic vocalizations, estrous staging, or a mating assay. What the result does establish is that a stable directional social dynamic exists in this behavior, and that characterizing it precisely benefitted greatly from multi-camera 3D pose estimation. 

Easily executable 3D pose methods are critical for enabling additional downstream methods such as embodied biomechanical modeling \citep{Leonardis2025,Zhang2025}. LUC3D successfully shifts the practical bottleneck by reducing manual labor and upfront cost of entry allowing for more users to begin their journey towards 3D pose tracking of social behavior. 

\section{Online Methods}\label{sec4}

\subsection{The application}\label{methods-app}

LUC3D is a multi-view pose annotation and 3D proofreading interface that runs entirely in a web browser. It is pure vanilla JS served as static files with no installation and no build step, either from the hosted deployment or from any local static file server, and it keeps no state on a server, so a project is simply a set of local files that the user opens and saves.

It loads multi-camera video with frame-accurate WebCodecs decoding and scrubbing synchronized across every view, camera calibration in TOML or JSON, and SLEAP predictions, which are read lazily so that a prediction file larger than available memory can still be opened. Several recording sessions can be held open at once, each in its own dockable pane. A project can also be opened with no video at all, as a skeleton together with imported 3D points, in which case the whole duration of the reconstruction remains navigable.

Skeleton instances and their keypoints are created and edited directly on the video frames, and the node and edge definitions themselves can be edited inline. The reprojection-aided labeling protocol quantified in Figure~\ref{fig2} is available in the interface: the user labels two views, the 3D point is solved from those two alone, and the reprojection appears in every remaining view to be accepted or nudged.

A cross-view tracker groups the detections that belong to the same individual across views and assigns one identity per animal per frame, which is the capability measured in Figures~\ref{fig3} and~\ref{fig6}. Its behavior is exposed through a guided dialog that offers per-node weights, per-view inclusion, and the cost thresholds described above, so that the association can be tuned to a particular rig rather than accepted as a fixed policy. Tracks and identities are held per session and propagate to one another, so that a correction to either is reflected in the other.

Triangulation is offered as the linear solve and as the non-linear refinement described above, with optional rejection of an outlying view, and the resulting reconstruction is drawn in an interactive Three.js viewport that also renders the calibrated camera frustums, can be placed at any camera's own pose and field of view so that it can be compared directly against that camera's video, and can be exported as an MP4 video. The proofreader's working signal is the reprojection error, surfaced for each keypoint and each view as an overlay on the frame being edited, with detections drawn as solid marks and reprojections from the 3D as dotted ones. The reprojection errors are also listed in an Instance info panel on the right side of the screen (See Supplementary Figure~\ref{fig7} for GUI screenshot). A timeline widget in the style of SLEAP carries frames, tracks and labeled regions for navigation, and the keyboard shortcuts are listed and can be rebound in the settings dialogue.

Three instance types are distinguished throughout the interface: user instances, created or edited directly by the annotator and drawn as solid, fully editable skeletons; predicted instances, produced by the cross-view tracker or imported from SLEAP and drawn dashed and read-only until converted; and reprojected instances, computed by projecting the current 3D reconstruction back into a view and drawn dotted and read-only. A new instance is placed at the view center of the camera under edit with a single key press, after which its keypoints are dragged into position individually, or, holding a modifier, moved together as a rigid group to reposition an already roughly correct skeleton. Each node can also be cycled through visible, occluded, and absent states with a right-click rather than being hidden off-frame, so that a keypoint's absence from a given view is recorded explicitly. Deletion is scoped to match this per-view structure: one shortcut removes the selected instance from the camera view under edit only, and a modified form of the same shortcut removes it from every view in its group at once, so a single mislabeled detection can be cleared without disturbing the same individual's instances in the other cameras.

The skeleton definition is edited inline from its own panel rather than as a separate file: nodes can be added, renamed, or removed, and edges rewired, with the change applied immediately to every open view. A project holds one skeleton shared across all of its sessions, so a definition entered for one camera rig carries over to the next session opened in the same project, and the interface warns before an edit, or a newly loaded skeleton, would overwrite one that already has instances attached to it.

The panel driving the reprojection-aided protocol reports, for every instance group, the reprojection error achieved in each contributing camera. A reprojected instance appearing in a still-unlabeled view is converted to an editable one with a double-click and nudged into place rather than positioned from scratch, after which the group can be re-triangulated with the added view to sharpen every remaining reprojection in turn. The same overlay is the proofreader's principal diagnostic: on an already-labeled, triangulated frame, a solid annotation that diverges visibly from its dotted reprojection flags either a mislabeled keypoint, a left–right identity swap, or, where the divergence is uniform across every view at once, a calibration that needs correction, and instances above a chosen error threshold can be reviewed frame by frame across an entire tracked session in a single pass.

The 3D viewport can also be rendered out rather than only inspected live: an export dialog offers a choice of output resolution from 360p up to 2K, reports the resulting frame count and an estimated file size before rendering begins, and encodes the sequence with WebCodecs directly in the browser, requiring only that the frames being rendered already carry triangulated 3D points or an imported reconstruction.

Projects are written as SLEAP .slp files, including the columnar session data introduced in version 2.8 of that format, which carries the calibration, the cross-view identities and the 3D points, so that a project round-trips through SLEAP and through this application without loss. Labels can also be exported as JSON, 3D points as HDF5, and the 3D viewport as MP4.

\subsection{Software}\label{methods-software}

The application is written as vanilla ES modules in JavaScript with no build step. It reads and writes version 2.8 of the SLP format through sleap-io.js and HDF5 through h5wasm, and renders in three dimensions with Three.js. A recent Chromium-based browser is recommended, because frame-accurate decoding uses WebCodecs. The benchmark drivers run the unmodified application modules under Node 26, so the association and triangulation measured here are the shipped code paths rather than a reimplementation of them. Analysis and figures were produced in Python 3.12 with NumPy, pandas, h5py, SciPy and Matplotlib. Tracking metrics come from motmetrics and the comparison triangulator from aniposelib 0.7.2.

\subsection{Datasets}\label{methods-datasets}

Three multi-camera recordings were used. No animal procedures were performed for this work, and all three recordings are pre-existing datasets that had already been proofread in 3D.

The Mouse-Dyad-10M corpus consists of 56 sessions, each recording two mice with five synchronized cameras at 150 frames per second for approximately 20 minutes, giving 10,084,734 frames in total. A 15-node mouse skeleton was used throughout, comprising the nose, left and right ears, the tail-tip and tail base (denoted TTI), three intermediate tail nodes, the head, trunk and neck, and the left and right shoulders and haunches. Figures~\ref{fig2},~\ref{fig3}E,~\ref{fig4} and~\ref{fig6}A were measured on this corpus.

This paper presents the SLAP-2M data repository of 2.81 million 3D poses and 22.49 million 2D poses generated by SLEAP-Anipose. The multi dataset was collected during 74 recording sessions with 8 camera views with between 1--4 animals in each session. The 8 camera views (back, mid, side, top, back-left, mid-left, side-left, and top) were pointed at the center of the environment. Both males and females are included in the dataset as well as multiple mouse strains (white, black, and agouti). This dataset is unique in its level of diversity and emphasis on freely roaming behavior. The dataset includes 6 different levels of enrichment: no enrichment, interactive planar enrichment, single animal and single vertical enrichment, multi-animal and single vertical enrichment, multi-animal and multiple vertical enrichment, multi-animal and multiple vertical and opaque enrichment. While many popular 2D datasets lack environmental enrichment, adding levels of enrichment highlights the need for 3D pose estimation. By adding more aspects of the environment which the animals can climb, 3D pose becomes more relevant. This makes the dataset an ideal candidate for benchmarking new models for behavioral clustering (both supervised and unsupervised), 3D pose estimation, and agent-based reinforcement learning algorithms. Each session also carries a curator-assigned difficulty rating and a black or white bedding label. The difficulty rating is the stratifying axis throughout Figure~\ref{fig5}, and the bedding split is reported for each stratum in Figure~\ref{fig5}F. Figures~\ref{fig5} and~\ref{fig6}B to F were measured on this corpus.

The third recording, HardFight, is a single eight-camera, three-mouse session at 60 frames per second, from which a 300-frame window was taken for Figure~\ref{fig1} and Figure~\ref{fig2}A. It is the only recording that was built as a complete application session with its own calibration, which is why the panels driven through the application itself use it.

\subsection{Animal Ethics Statement}\label{methods-mice}

All animal procedures were approved by the Princeton University Institutional Animal Care and Use Committee in accordance with National Institute of Health standards.

% TODO: mouse strains/conditions/behaviors, length and size of video data, storage, session structure.

\subsection{Camera calibration}\label{methods-calibration}

Calibration is the process by which camera parameters and 3D points are iteratively estimated by comparing reprojected 2D points with multi-view ground truth. The ground truth is estimated by using an easily detected stimulus, a checker board with bit encodings called Aruco patterns. A calibration is performed by pointing multiple cameras at a central point, and the ChAruco board is presented in multiple orientations and positions visible from all cameras. Then sparse bundle adjustment is used to estimate optimal camera parameters and 3D points which are associated with unique positions of the detected markers. The camera calibration is stored in TOML format which records the camera name, resolution, intrinsics, extrinsics (rotation and translation), and distortion parameters.

Each rig ships a calibration file that gives, for every camera, its intrinsic matrix, its radial and tangential distortion coefficients and its extrinsic pose, in a TOML format compatible with aniposelib. The application reads these files unchanged, and no recalibration was performed for this work. Reprojection errors are reported throughout in each camera's own native, still-distorted pixels, because that is the space in which an annotator sees the error; a point is undistorted into ideal-pinhole coordinates only inside the linear solver, and the residual is measured after it is projected back. All 3D coordinates are expressed in the calibration's own metric frame, in millimeters. For Figure~\ref{fig2}C the reconstructed pose was aligned to the calibration frame by a RANSAC-Procrustes fit, which retained 98.3 per cent of correspondences as inliers and left a residual of 1.32~mm. This paper presents a browser-based calibration tool which takes OpenCV and Anipose-style calibration suite and makes it more accessible without installation at \url{https://vibes.tlab.sh/calibration-studio/}.

\subsection{Two-dimensional detection pool}\label{methods-pool}

Every benchmark reported here ran on a single shared detection pool from which all track and identity information had been removed, so that the comparisons isolate cross-view association and triangulation from the quality of the underlying detector. Each method under comparison received byte-identical input. Ground-truth identities were transferred onto that pool by matching detections to the proofread instances frame by frame and camera by camera at an intersection-over-union threshold of 0.5, so that a detector identity swap appears in the results as an error of the method under test rather than being silently absorbed. SLEAP produced the detections; the application consumes SLEAP .slp files and neither trains nor runs a detector of its own.

The measurements divide into two kinds, and they treat the recordings differently. The association and tracking results in Figures~\ref{fig3} and~\ref{fig6}, and the behavioral analysis in Figure~\ref{fig4}, were computed on complete sessions, every frame, with no subsampling. The per-keypoint geometric measurements in Figures~\ref{fig2},~\ref{fig4} and~\ref{fig5} were computed on a uniform sample of frames, every 15th frame for Figure~\ref{fig4}; Figures~\ref{fig2} and~\ref{fig5} use every frame, except that Figure~\ref{fig2}C is drawn from Figure~\ref{fig4}'s by-camera-count measurement pass (the $k$-view solves and their held-out scores) and therefore carries Figure~\ref{fig4}'s sampling and its 50 proofread sessions. The sample is uniform and was never selected on any property of the data, and it is large in absolute terms: 286,200,174 keypoints in Figure~\ref{fig2}, 17,013,412 in Figure~\ref{fig4}, and 187,134,382 keypoint comparisons in Figure~\ref{fig5}. Figure~\ref{fig4}'s subsampling was necessary because these measurements are per-keypoint and combinatorial in the cameras: Figure~\ref{fig4}B rests on 884,697,424 solves at every 15th frame, and an every-frame pass would multiply that count by fifteen. The every-frame measurement in Figure~\ref{fig2} rests on 2,862,001,740 two-anchor solves and 8,586,005,220 held-out view measurements.

The sampled density was verified rather than assumed. The 3D error by anchor count computed for Figure~\ref{fig4}B at every 15th frame reproduces the same quantity computed for Figure~\ref{fig2} at every frame to within 0.13 per cent at every camera count, comparing medians of per-session medians over the same 50 sessions. This is the expected result for these recordings, since at 150 frames per second an animal moves well under a millimeter between consecutive frames, so a stride of 15 samples the behavior every 0.1~s and discards redundancy rather than information. The one place where the stride interacts with a method rather than merely with the sample size is temporal smoothing, and that is the reason the smoothing term is disabled in the comparison triangulator, as described below.

\subsection{Cross-view association}\label{methods-association}

For each frame the application maintains a set of 3D targets and solves one Hungarian assignment per camera, committing each camera's assignment before the next is solved, at a cost of $O(C \cdot A^3)$ for $A$ animals and $C$ cameras. The cost of pairing a target $t$ with a detection $d$ is summed over the skeleton nodes $k$ with per-node weights $w_k$, and negated so that the solver minimizes it. The two-dimensional term compares the detection with the target's projection $\pi(t)$ into that camera, in undistorted ideal-pinhole coordinates,
\[
  w_k \cdot corr2d \cdot \left(1 - \frac{|d_k - \pi(t)_k|}{velThresh \cdot (1 + \Delta t)}\right) \cdot e^{-\lambda \Delta t}
\]
where $\Delta t$ is the target's age and $\lambda$ the time penalty, so age both widens the distance allowance and decays the whole term. The three-dimensional term is on the perpendicular distance from the target's 3D point to the ray back-projected from the detection, with no age decay,
\[
  w_k \cdot corr3d \cdot \left(1 - \frac{dist(t_k, ray(d_k))}{distThresh}\right)
\]
Both thresholds are soft, so exceeding one drives its term negative rather than preventing the pairing. The ratio $r = corr3d/corr2d$ controls how strongly the association relies on 3D ray geometry relative to 2D reprojection distance. Two further properties complete the association. A retained per-view detection older than 20 frames is evicted from a target's 3D anchor before that frame's matching runs, so an anchor is only ever triangulated from cameras that saw the animal recently; and every camera is scored against one frame-start snapshot of the anchor, which is re-triangulated once at frame end, so a decision in one view cannot move the evidence the remaining views are judged against. The application ships $corr2d = 1$, $corr3d = 6$, velocityThreshold 10, $distThresh = 25$, timePenalty 0.1 and the 20-frame eviction window, and every measurement in this paper uses them.
 
The method compared against, hereafter exhaustive enumeration, is our reimplementation of the per-frame procedure of \citep{Maree2024}. It enumerates all $(A!)^C$ groupings of a frame's detections into identities, triangulates and reprojects each one, and keeps the grouping with the lowest error. It is undefined unless every camera holds exactly $A$ detections, so only such frames were eligible and ineligible frames were counted rather than dropped. The benchmark covers the 92 multi-animal sessions carrying both pool detections and proofread ground truth, 50 Mouse-Dyad-10M and 42 SLAP-2M. Of their 9,678,503 frames, 4,591,864 were eligible and all were computed, the three- and four-animal configurations at 2.7 and 5.5 seconds per frame against 11 to 16 milliseconds for two animals. Four animals in six cameras exceeds the harness cap of $10^6$ hypotheses per frame by a factor of 191 and was not run, so its cost in Figure~\ref{fig3}F is an arithmetic bound drawn with an open marker, from the $(A!)^{(C-1)} = 7{,}962{,}624$ hypotheses remaining once the global relabeling symmetry is removed up to the as-published $(A!)^C$.
 
Grouping quality in Figure~\ref{fig3}E compares partitions rather than labels, so a consistent relabeling of identities is not counted as an error, and the 4,591,725 computed frames carrying a transferred ground-truth match are scored. Two conventions separate the arms. First, exhaustive enumeration has no cross-frame identity of its own, so the harness threads identity between consecutive frames by nearest-3D-centroid matching. That threading is scaffolding rather than part of the published method, leaving the per-frame agreement rate as the like-for-like comparison. Second, the arms are scored over different exposures, exhaustive enumeration over the 21,622,345 clean camera-frames it computed and the greedy arms over whole sessions, 45,021,960 camera-frames, so each rate carries its own denominator.
 
The sweep in Figures~\ref{fig3}G and H fixes $corr2d = 1$ and varies $corr3d$ to sample twelve values of $r$ over all 50 Mouse-Dyad-10M sessions at full length. Only the ratio matters, not the two weights individually, verified on an earlier 24-cell grid over eight sessions where every pair of cells sharing a ratio returned identical IDF1 and identical switch counts. Switches are reported per 100,000 camera-frames, with the denominator taken per camera and session from the scorer's own frame counts and checked identical between the arms before they are drawn on one axis.

\subsubsection{Camera-subset identity}\label{methods-subset}

The number of cameras identity needs was measured by running tracker on camera subsets of the Mouse-Dyad-10M rig: $k = 2, 3$ and 4 cameras with three fixed subsets per $k$, plus $k = 5$, the full rig, over the same 50 sessions and the same shared detections. The subsets are deterministic: the $C$-choose-$k$ combinations of the five camera serials are sorted and taken at the first, middle and last index, with no random draw, so the run is reproducible from the script alone, and the same three subsets are used for every session, so between-session spread is not confounded with between-subset spread (the full 26-subset design at 50 sessions, 1,300 runs, was ruled infeasible). The $k = 5$ cell is not re-tracked: it is required to reproduce default tracker's deposited 50-session numbers exactly, and does (maximum absolute IDF1 difference 0.0 against the reference deposit), which is what makes the subset cells attributable to the subset rather than to harness drift. Within-view IDF1 rises with the cameras available: subset means 0.669 to 0.688 at $k = 2$, 0.726 to 0.743 at $k = 3$ and 0.730 to 0.776 at $k = 4$, against 0.749 for the full rig (per-$k$ means of roughly 0.68, 0.74, 0.75 and 0.75). Because the exposure shrinks with $k$, switch counts from different $k$ are compared only as rates over the camera-frames of the cameras used, never as raw sums. The panel drawn from this measurement is not yet placed on any figure.

\subsection{Triangulation}\label{methods-triangulation}

The application's default linear solver undistorts the observations into ideal pinhole camera coordinates and minimizes the algebraic error in closed form by direct linear transformation. The non-linear solver minimizes a geometric error in each camera's native, still-distorted pixels while holding the cameras fixed, using a soft-L1 loss followed by an L1 polish phase and initialized from the linear solution, with a backtracking guard that vetoes any step which raises the objective. Although the application's menu labels it ``Bundle Adjustment,'' the cameras are frozen so is more akin to a refined iterative triangulation. 

The application can optionally drop the view with the worst residual and re-solve the point from the remaining views. Figure~\ref{fig2}E scores the same solve before and after dropping the camera with the least reliable detections and shows an overall decrease in reprojection error. 

The comparison triangulator in Figures~\ref{fig2}F and~\ref{fig4}G is aniposelib 0.7.2, the OpenCV and NumPy release that undistorts with \texttt{cv2.undistortPoints} and then solves each point with a singular value decomposition, which is also CPU based. While there are also GPU-based alternatives those would not be suitable for comparison due to the browser-based constraint. The four solvers were compared in pairs matched by algorithm class by linear and non-linear forms. Enabling the library's RANSAC path costs 2,467 microseconds per keypoint, which is 85 times the default path, and it is reported here rather than drawn because a bar of that size would flatten every other bar in the panel. The temporal smoothing term in the non-linear path was disabled, because the smoothing penalty treats consecutive entries as consecutive video frames, and the input is sampled every 15th frame and filtered to keypoints complete in all five views, so consecutive entries are separated by 0.1~s of real motion or more and the term would penalize that motion as though it were noise; with the smoothing left enabled the library's optimizer is worse than its own linear solve in all 50 sessions. The rule was fixed before the results were seen: the smoothing term is invalid on strided input and is disabled. Both variants were measured and deposited, so the choice can be checked.

Solve times exclude undistortion for every solver, because one library undistorts inside the triangulation call and the other outside it, and charging only one of them would measure where each library draws a function boundary rather than the cost of the solve; the excluded undistortion is 0.57 microseconds per keypoint for aniposelib and 1.16 for this work. The solvers reported here were timed in process under single-threaded Node 26 with \texttt{performance.now} around each call, and aniposelib with \texttt{time.perf\_counter}, taking the best of three runs. The two libraries are written in different languages and run on different runtimes, so this comparison is of two implementations as a user would encounter them rather than of two algorithms; the browser deployment that motivates this work forecloses a compiled implementation, which is why the shipped solver is the one measured.

Figure~\ref{fig2}C, by contrast, reports the 3D distance to the proofread reference in millimeters, deliberately: there the comparison is across camera counts within one solver rather than between solvers, so the objection above does not apply, and the reference's own noise (its median reprojection error is 2.40~px) enters every camera count equally. Its absolute level is therefore a comparison value rather than absolute 3D accuracy, which its legend states.

\subsection{Reprojection recovery}\label{methods-recovery}

Figure~\ref{fig5}C measures how many of the detector's per-view misses the rest of the rig could fill in by reprojection, as a function of rig size. For every keypoint-instance (one node of one animal in one frame), $g$ is the number of cameras whose proofread reference carries that keypoint and $m$ is the number of those cameras whose matched detection carries it too. The matching is identical to the Figure~\ref{fig5} detection-quality convention, a mean-keypoint-distance Hungarian assignment against the proofread reference, run at stride 1, so ``missing'' in this measurement is the same notion the detection-quality panels plot. Each session deposits its $(g, m)$ histogram, and the histograms are the whole measurement: every rig-size statistic is exact arithmetic on them, with no sampling and no per-subset re-runs. A view that misses a keypoint counts as recovered when at least two of the other cameras in the rig detect it, because the keypoint can then be triangulated from those views and reprojected into the view that missed it; one other view cannot triangulate, so recovery at a rig of two cameras is zero by construction. For a rig of $k$ of the six cameras, the probability that a missing view is recoverable is an exact hypergeometric expectation over all $C(6,k)$ camera subsets: conditioned on the missing view being in the subset, the other $k-1$ slots draw from the remaining cameras, of which $m$ detect, and recovery requires at least two detecting among them. Pooled over the 74 sessions, 42,184,875 keypoint-instances, 26.1 per cent of keypoints are missing per view; recovery rises from zero at $k = 2$ to 45.7 per cent of misses at $k = 6$, leaving 14.2 per cent of keypoints still missing at the full proofread rig.

\subsection{Behavioral analysis of social rearing courtship behavior.}\label{methods-behavior}

A frame was counted as a rear when the animal's neck height exceeded 0.75 of its own body length, where body length is that animal's median nose-to-tail-base distance across the session. Rearing bouts shorter than 0.25~s were discarded, and gaps of 0.15~s or less within a bout were merged. A mutual upright display is both animals reared with their tail bases within two body lengths of each other, held for at least 0.25~s. There are 539 such displays in 37 of the 56 Mouse-Dyad-10M sessions.

Because a display begins only when the second animal comes up, the onset of the display itself cannot identify which animal started it. The initiator was therefore identified from each animal's own rearing bout containing that onset, taking whichever animal's bout began earlier, and the lag reported in Figure~\ref{fig4}C is the interval between the two bout onsets.

A session's leader, in Figure~\ref{fig4}F, is whichever of the two animals initiated more of that session's displays. Because that statistic is the larger of the two shares it cannot fall below 0.5, it is read against a size-matched surrogate rather than against a line at 0.5. Each surrogate corpus keeps every session's own display count and only relabels which animal started each display, drawing from a binomial distribution with probability one half and taking the larger share exactly as the observed statistic does, so the surrogate carries the same floor and the same small-session inflation as the data. Two thousand surrogate corpora were drawn and none reached the observed median, which is the basis for the bound of $P < 0.0005$ quoted in the caption. The figure reported in the text is pooled over displays; averaging over sessions would weight every session equally regardless of its display count. The per-session statistics are restricted to sessions containing at least six displays, which is the smallest session in which a two-sided binomial test can reach significance at all, since the best attainable $P$ value is 0.0625 at five displays and 0.031 at six. That leaves 23 sessions. The leader's share is individually above chance by a two-sided binomial test in 16 of those 23 sessions, and in 9 of 23 after Holm correction across the family.

The 37 sessions are repeated recordings of a smaller number of pairs, so the same analysis was repeated with the pair as the unit of replication. Aggregating every session of a pair into one observation gives 14 pairs with at least five displays, covering 536 displays. The leading member of the pair starts a median 0.81 of that pair's displays (interquartile range 0.74 to 0.91, range 0.60 to 1.00), and the same member leads in all 14 of 14 pairs (sign test $P = 1.2 \times 10^{-4}$; Wilcoxon signed-rank on the pair shares against 0.5, $P = 9.7 \times 10^{-4}$). Eight of the 14 pairs are individually above chance by a two-sided binomial test, and all eight remain so after Holm correction. Pooled over those pairs the leader starts 429 of 536 displays, 80.0 per cent, which is the same figure the session-level pooling gives. The result is therefore not an artifact of treating repeated recordings of one pair as independent.


The coupling in Figure~\ref{fig4}G was computed by taking every rearing onset by one animal and reading out the probability that the other animal was rearing at each lag within five seconds either side, divided by that other animal's own base rate. Onsets were split by the tail-base separation at the moment of onset, into those within two body lengths and those beyond. The null distribution was generated by circularly shifting the other animal's rearing time series, 24 shifts per ordered pair, which preserves that animal's rearing rate, its bout durations and its autocorrelation while destroying only the temporal alignment between the two animals; a reshuffle would additionally have destroyed the autocorrelation and would have made almost any structure appear significant. A session contributed a curve to a condition only if it supplied at least 20 onsets in that condition. The same measurement on the two-animal sessions of SLAP-2M gives 1.08 within two body lengths and 0.97 beyond, but that arena is 3.2 body lengths across against Mouse-Dyad-10M's 6.9, so the two conditions barely differ there and the claim is made for Mouse-Dyad-10M alone.

\subsection{Tracking metrics}\label{methods-tracking}

IDF1, identity switches and fragmentations were computed with the motmetrics library on the shared detection pool. Within-view IDF1 was computed separately for each camera and averaged. Cross-view IDF1 pools all cameras into a single accumulator with one global identity per animal. For a tracker that has no cross-view identity of its own, that pooling is camera-scoped: its hypothesis ids are keyed by camera, so the same track number appearing in two cameras is two identities, and global ids are used only for methods that actually assert cross-view identity. The scoped convention is the conservative one: unscoped pooling can only add cross-view matches, and with two enforced tracks per camera it would credit the capped SLEAP baseline of Figure~\ref{fig6}A with a cross-view IDF1 of 0.600 through nothing but shared slot numbering. Under scoped pooling a method that labels each camera independently can be matched to the truth in at most one camera, which puts such a method near $1/C$; that level is a property of the pooling convention rather than a strict bound, and chance depends on the number of animals alone. A fragmentation was counted each time a tracked ground-truth track became untracked and was later picked up again, which is a different event from an identity switch, in which the track continues but is assigned to the wrong animal. Error terms are reported as percentages of camera-frames. False negatives were measured but are not plotted, since they account for 98.8 to 99.3 per cent of every method's error budget and would draw three indistinguishable bars.

ID accuracy (IDA) is idtp divided by the number of matches: of the detections matched to a ground-truth animal, the fraction carrying the correct identity. It excludes both detector misses, which sit in IDR's denominator, and false-positive detections, which sit in IDP's, so it isolates the association from the detector on both sides, and IDA $\geq$ IDP always. For the fresh-anchor configuration on the 50 Mouse-Dyad-10M sessions, pooled IDA is 92.46 per cent (72,794,704 of 78,734,134 matches), the session median is 100 per cent and the worst session 56.6 per cent; false-positive detections are 0.104 per cent of matches (81,754).

\subsection{Baseline configuration}\label{methods-baseline}

The per-camera baselines in Figure~\ref{fig6}A were re-run at the most favorable configuration available to them, because the comparison as first run charged them for constraints LUC3D's arm does not face: LUC3D holds two global identities by construction, while SLEAP ran with an unbounded track pool and ByteTrack with a 2-second retirement horizon on 20-minute sessions.

SLEAP was re-run with the track count capped at the true animal count. The original retrack ran sleap-nn 0.3.0 with no cap (its \texttt{--max\_tracks} option is honored only under \texttt{candidates\_method} `local\_queues', and the default `fixed\_window' silently ignores it), yielding a median of 47 tracks per camera-session. The re-run adds exactly two flags, \texttt{--max\_tracks 2 --candidates\_method local\_queues}, and changes nothing else, so the only difference between the runs is the cap; the cap was verified after the fact, with all 250 camera-sessions holding at most 2 tracks. The output is scored from the tracked .slp files directly, every track.

ByteTrack was re-run with track retirement disabled, setting \texttt{lost\_track\_buffer} to the session length against the shipped 60 frames (2~s), which retires a track on any longer occlusion, and its output was then reduced to two identities by a greedy, tracklet-whole stitch that uses no ground truth: each ByteTrack id is bound at birth to one of two slots, the slot whose last-seen box is nearest the tracklet's first box (intersection-over-union first, centroid distance as the tie-break), among slots not held by an id alive in the same frame, and keeps that slot for life, so ByteTrack's own association is untouched. The never-retire knob alone reaches within-view IDF1 0.272; with the stitch, 0.676. Two gates anchor the re-runs to the shipped measurement: re-scoring the shipped ByteTrack arm through the re-run's scorer reproduces the reference evaluation to a maximum absolute difference of 0.0 on all 50 sessions, and re-running one camera-session at the shipped parameters under the pinned library version reproduces the stored output cell for cell.

\subsection{Anchor staleness and the fresh-anchor tracker}\label{methods-anchor}

The cross-view tracker follows the tracking-by-3D-consensus design of \citep{Chen2020}. The browser port as first written retained, for every identity and every camera, the last detection matched to that identity, and re-triangulated the identity's 3D anchor from all retained detections with no age limit. Chen et al (2020) \citep{Chen2020} instead attenuate each retained point's contribution to the incremental triangulation exponentially with its age (their Eq.~11), but the reference implementation computes those weights without applying them, so retained detections of any age entered the anchor at full strength, and the browser port reproduced that behavior. We note additionally that the reference computes the weight with inverted sign, so that older points would receive larger weights were the routine completed as written. Logging the age of every detection entering a re-triangulation over eight Mouse-Dyad-10M sessions gave per-session mean ages of 3.0 to 49.8 frames and maxima of 844 to 8,652 frames, about five minutes of recording. Three coordinated changes were evaluated against that defect. Staleness eviction deletes, at the start of each frame, every retained detection older than $N$ frames, with $N$ swept over 1, 10, 20 and 30, so an anchor can be triangulated only from cameras that saw the animal recently, which is the step-function limit of the exponential attenuation of 3D correspondence; a target left with fewer than two retained detections keeps its last anchor but stops being re-triangulated, and an animal lost outright re-enters through the epipolar initialization path rather than by bidding from an obsolete position. Synchronous scoring, written M1, snapshots every anchor at frame start, scores all views against the snapshot and re-triangulates once at frame end, so a wrong assignment in one view cannot contaminate the evidence the remaining views are scored against and the result no longer depends on camera order. The distance threshold of the 3D term was reduced from 50 to 25, tightening the acceptance region around any anchor that does freeze. The assignment algorithm itself was unchanged, so every effect arises from the evidence entering the affinity computation. The evicting configuration these runs measured through that loader hook is the one the application now ships, merged into \texttt{pose/cross-view-tracker.js} itself, so every measurement reported in this paper describes the released tracker. All configurations were run on the 50 proofread Mouse-Dyad-10M sessions with identical per-camera detections, in a harness that overlays each candidate on the production module through a loader hook; before any comparison the harness reproduced the unevicted baseline's deposited 50-session result bit-identically. Paired comparisons across the 50 sessions used the two-sided Wilcoxon signed-rank test.

\subsection{Statistics}\label{methods-stats}

The session is the unit of analysis throughout unless stated otherwise. For Mouse-Dyad-10M this is a conservative-sounding choice that is not automatically conservative, because the 56 sessions are repeated recordings of 9 mice in 18 pairings; where a session-level test carries a claim, the same test with the pair as the unit is reported beside it (Figure~\ref{fig4}F). For SLAP-2M the individual animals are not identified in the metadata, so no finer unit is available. Medians and interquartile ranges are used for the error distributions, which are right-skewed, and a band labeled as running from the 25th to the 75th percentile is an across-session spread of per-session medians rather than a confidence interval. The paired comparisons in Figures~\ref{fig6}A,~\ref{fig6}C and~\ref{fig6}F carry bootstrap confidence intervals at 95 per cent over sessions; the paired per-session comparisons of direction in Figure~\ref{fig6}C use sign tests; and the paired session-level comparison against a fixed value in Figure~\ref{fig4}F uses a Wilcoxon signed-rank test. Binomial nulls were simulated rather than approximated wherever the statistic is a larger-of-two share, using 20,000 draws per session size, and the temporal coupling in Figure~\ref{fig4}G is read against circular-shift nulls. No correction for multiple comparisons is applied across figures, and where a family of per-session tests is reported both the uncorrected and the Holm-corrected counts are given.

\backmatter

\bmhead{Acknowledgments}
% TODO: add acknowledgments.

\bmhead{Data Availability}
% TODO: add data availability statement (repository links for Mouse-Dyad-10M-11M, SLAP-2M, HardFight).
The pose estimation data from the SLAP-2M dataset is available at \hyperlink{https://huggingface.co/datasets/talmolab/slap-2m}{https://huggingface.co/datasets/talmolab/slap-2m}. The pose estimation data from the Mouse-Dyad-10M dataset is available at \hyperlink{https://huggingface.co/datasets/talmolab/mouse-dyad-10M}{https://huggingface.co/datasets/talmolab/mouse-dyad-10M}.

\bmhead{Code Availability}
% TODO: add code availability statement (LUC3D repository link).
The luc3d GUI can be accessed directly at \hyperlink{https://talmolab.github.io/luc3d/}{https://talmolab.github.io/luc3d/}. The code for the luc3d can be found at \hyperlink{https://github.com/talmolab/luc3d}{https://github.com/talmolab/luc3d} with documentation at \hyperlink{https://talmolab.github.io/luc3d-docs/}{https://talmolab.github.io/luc3d-docs/}. 

\bmhead{Supplementary Movies}

Available online at: \url{https://talmolab.github.io/luc3d-page/}

% ---- Supplementary material -------------------------------------------------
\section{Supplementary Material}\label{sec5}

\subsection{Cross-view identity tracking: parity within cameras, gains across them}\label{subsec2-7}

Supplementary Figure~\ref{fig6} compares LUC3D against per-camera trackers on identical, identity-stripped detections, so that what is compared is association rather than detection. The cross-view comparison of Figure~\ref{fig6}A is measured on the 50 Mouse-Dyad-10M sessions, a fixed five-camera rig with two animals throughout; the per-session comparisons of Figure~\ref{fig6}B to F are measured on all 74 SLAP-2M sessions, which run one to four animals over eight cameras of which six are proofread and span the full curator difficulty range. Absolute IDF1 is accordingly lower for every method on SLAP-2M than on Mouse-Dyad-10M, and the two corpora's numbers are not comparable to each other.

Within a single camera of SLAP-2M, all three methods are broadly competitive, and LUC3D's advantage over SLEAP is largest in the single-animal sessions where there is nothing to associate (Figure~\ref{fig6}C). Pooled over all 74 sessions the difference in within-view IDF1 is 0.099 in LUC3D's favor (56 of 74 sessions, sign test $P = 1.1 \times 10^{-5}$). Restricted to sessions with two or more animals it is 0.067, ahead in 31 of 42 sessions ($P = 0.003$), and in the two cells where cross-view association ought to help most, three and four animals, it is still negative, with SLEAP ahead in all seven of those sessions, though the deficit roughly halves against the unevicted tracker (medians $-0.039$ and $-0.028$). We state that rather than pooling it away.

The difference appears when identities are pooled across cameras (Figure~\ref{fig6}A). Scoring all five Mouse-Dyad-10M cameras into one accumulator with a single global identity per animal, LUC3D's IDF1 is unchanged, 0.861 within view and 0.861 across at the fresh-anchor configuration the panel draws. The per-camera trackers lose about three quarters or more of theirs, 0.642 to 0.146 and 0.676 to 0.157, a retention of 0.23. This is close to arithmetic rather than a contest, and we do not present it as a benchmark win: a tracker that labels each camera independently can have its labeling matched to the truth in at most one camera, so $1/C$ is a ceiling for such a method, and the comparison scores those methods on a capability they never claimed. What the comparison establishes is that the capability is absent; it does not show the methods to be poor. Per-camera identity and cross-view identity are different objects, and only the latter supports a 3D reconstruction or a proofreading pass that acts on all views at once.

There is a cost. LUC3D fragments tracks more than SLEAP does, by a mean of 3.2 fragmentations per camera-session and a median of 0.9, and SLEAP fragments less in 70 of the 74 sessions. A fragmentation differs from an identity switch: the track breaks and is later resumed rather than being handed to the wrong animal. For a proofreading tool that is the preferable failure, since a break is visible and locally repairable while a switch propagates silently, but it is a real cost and it is the price of the same conservatism that produces the cross-view result.

\subsection{Stale evidence in the 3D anchor is the dominant switch mechanism}\label{subsec2-9}

The identity errors that dominate the tracker's remaining deficit have a single cause, and it is visible in the parameter sweep of Supplementary Figure~\ref{fig6}D. Because the tracker as originally ported re-triangulated an identity's 3D anchor from retained detections of any age, and because the 3D affinity term carried no temporal attenuation, an identity that had lost its animal minutes earlier kept bidding at full strength from the last available position. When the other animal passes through that position the joint assignment concedes it, and with no appearance cue available the swap is introduced. Evicting retained detections older than $N$ frames before matching, together with synchronous per-frame scoring and the tightened 3D distance threshold, reduced within-view identity switches from 2,071 to 413, from 0.0046 to 0.0009 per cent of camera-frames, at $N = 20$, and raised mean cross-view IDF1 from 0.749 to 0.861 across the 50 sessions. At the shipped $N = 20$ the median session rises from 0.760 to 0.915; 38 of the 50 sessions improved and 8 worsened (two-sided paired Wilcoxon, $P = 1.7 \times 10^{-6}$), and the neighbouring window $N = 10$ gives 511 switches and mean IDF1 0.850, and the margin over the best pure re-tuning of the unevicted tracker was a median of 0.078 IDF1 ($P = 2.7 \times 10^{-5}$) (See Supplementary Figure~\ref{fig6}A-C).

Two features of the sweep argue that the effect is specific to anchor freshness rather than to any particular setting. The eviction windows are nearly indistinguishable from one another, and their order is not stable: $N = 1$ and $N = 10$ reverse between an eight-session pilot and the 50-session benchmark (0.839 against 0.850, $P = 0.8$), so the material step is from no eviction to any eviction rather than the choice of window. Traversing the same axis in the opposite direction is monotonically harmful across four orders of magnitude, with exponential smoothing of the anchor raising switches to roughly 2,800 and freezing the anchor on low-margin frames giving roughly 1,700 at margin 100 and 27,042 at margin 400. The application ships $N = 20$, which is best on every aggregate and on the harm profile alike: 413 switches against 511 at $N = 10$, mean cross-view IDF1 0.861 against 0.850 and median 0.915 against 0.913, and 8 rather than 10 of the 50 sessions worsened, its one cost being a deeper worst single-session change ($-0.183$ against $-0.138$). The mechanism corrects an implementation omission rather than introducing a new principle: the omitted age weighting is benign over the 0 to 300~ms that retained points span in the original setting, and dominant at behavioral time scales, where occlusions leave evidence stale for minutes.

\begin{figure}[!htb]
  \centering
  \includegraphics[height=.49\textheight]{figs/fig6.pdf}
  \caption{\textbf{Cross-view identity, and the fresh-anchor staleness horizon.} \textbf{A}, Within-view against cross-view IDF1, mean +- 95\% bootstrap CI over sessions: LUC3D 0.861 to 0.861 (stale 20, distThresh 25, synchronous scoring), SLEAP 0.642 to 0.146, ByteTrack 0.676 to 0.157. The dashed rule at 1/C = 0.20 is the camera-scoped pooling convention for a per-camera tracker, not a chance level. \textbf{B}, Within-view IDF1 as a survival curve, one step per session, each session's IDF1 the mean over its 6 cameras; at 0.9 or above, 40, 22 and 10 of 74 (LUC3D, SLEAP, ByteTrack). \textbf{C}, Paired within-view difference, LUC3D minus SLEAP, by animal count. Boxes, the per-session differences (median, IQR, whiskers 1.5x IQR; every session also a dot): medians +0.143, +0.068, -0.039, -0.028 at 1 to 4 animals (paired means +0.142, +0.089, -0.043, -0.041), LUC3D winning 25/32, 31/35, 0/4 and 0/3 sessions; over 2 animals, mean +0.067 (n = 42, 31/42, P = 0.003). \textbf{D}, False positives and ID switches as percentages of camera-frames: 0.317, 0.531, 0.435\% and 0.0112, 0.0308, 0.105\%. False negatives (98.8-99.3\% of each budget) are not plotted. \textbf{E}, The tracker's two association cost terms on one real frame (drawn in the style of Chen et al., 2020): the 2D term against the retained per-view 2D anchor (the dashed translucent blue outline pose), and the 3D term as the detection's back-projected ray against the 3D anchor node. The retained anchor is what panel F's staleness window evicts. \textbf{F}, The fresh-anchor parameter sweep, all 50 proofread Mouse-Dyad-10M sessions, 45,021,960 camera-frames, sessions paired across parameter sets; every candidate is M1 + distThresh 25 + the stale window named in the key. Left, the identity precision-recall plane: one operating point per parameter set at the across-session median (identity recall x, identity precision y), whiskers the IQR on both axes (operating points, not a swept curve: this detection pool has no score to threshold). Middle, cross-view IDF1 as a survival curve. Right, ID switches as a percentage of camera-frames; the raw totals behind the percentages are 2,071 (no eviction), 677, 511, 413 and 487 (stale 1, 10, 20, 30) -- no eviction 0.00460\% against stale 20 at 0.00092\%, 80\% lower, with median cross-view IDF1 0.7604 to 0.9153. A, 50 Mouse-Dyad-10M sessions; B to D, all 74 SLAP-2M sessions, 11,726,640 camera-frames; E is one frame of a two-animal SLAP-2M session; F, the 50 proofread Mouse-Dyad-10M sessions. The tracker key under panels A to D (LUC3D teal, SLEAP periwinkle, ByteTrack orange) names the three colours once for all four panels; F's parameter-set key is the strip along its own bottom.}
  \label{fig6}
\end{figure}

\begin{figure}[!htb]
  \centering
  \includegraphics[height=.4\textheight, width=1\textwidth,]{figs/fig7.png}
  \caption{Screenshot of the LUC3D GUI.}
  \label{fig7}
\end{figure}

% A table float would swallow this caption: sn-jnl only emits a table's stored
% \caption when a tabular begins inside the float, and this table is an image.
% So it floats as a figure and is captioned as a table via \captionof.
\begin{figure}[!htb]
  \centering
  % generated by figs/dataset_sheets.py (source: figs/figures/datasheets/datasheet_combined.pdf)
  \includegraphics[width=0.95\textwidth]{figs/datasheet_combined.pdf}
  \captionsetup{font=scriptsize}
  \captionof{table}{\textbf{Dataset datasheets.} Composition, metadata and per-session file inventory of the two corpora. SLAP-2M: 74 sessions, 8 camera views (2 held out from proofreading), 1--4 animals. Mouse-Dyad-10M: 56 sessions, 5 camera views, 2 animals, 10.08 million frames per view. 2D pose counts are reprojections of the 3D poses into every view.}
  \label{tab:datasheets}
\end{figure}

\FloatBarrier

\begin{thebibliography}{99}

\bibitem{Aldarondo2019} Aldarondo, D. (2019). Label3D. https://github.com/diegoaldarondo/Label3D

\bibitem{Aharon2026} Aharon, L., Whiteway, M. R., Sikka, K., Lee, K., Wang, Y., Chettih, S., et al. (2026). Lightning Pose 3D: an uncertainty-aware framework for data-efficient multi-view animal pose estimation. \textit{bioRxiv}.

\bibitem{Biderman2024} Biderman, D., Whiteway, M. R., Hurwitz, C., Greenspan, N., Lee, R. S., Vishnubhotla, A., et al. (2024). Lightning Pose: improved animal pose estimation via semi-supervised learning, Bayesian ensembling and cloud-native open-source tools. \textit{Nature Methods}, 21(7), 1316--1328.

\bibitem{Chen2020} Chen, L., Ai, H., Chen, R., Zhuang, Z., \& Liu, S. (2020). Cross-view tracking for multi-human 3D pose estimation at over 100 FPS. In \textit{Proceedings of the IEEE/CVF Conference on Computer Vision and Pattern Recognition (CVPR)}, 3276--3285.

\bibitem{Daruwalla2026} Daruwalla, K., Nozal Martin, I., Zhang, L., Naglič, D., Frankel, A., Rasgaitis, C., et al. (2026). Cheese3D enables sensitive detection and analysis of whole-face movement in mice. \textit{Nature Neuroscience}, 1--12.

\bibitem{Dias2025} Dias, I. C., Gutierrez-Castellanos, N., Lenschow, C., \& Lima, S. Q. (2025). Ready or not: Neural mechanisms regulating female sexual behavior. \textit{Current Opinion in Neurobiology}, 93, 103069.

\bibitem{Dunn2021} Dunn, T. W., et al. (2021). Geometric deep learning enables 3D kinematic profiling across species and environments. \textit{Nature Methods}, 18, 564--573.

\bibitem{Huang2020} Huang, C., Jiang, S., Li, Y., Zhang, Z., Traish, J., Deng, C., et al. (2020). End-to-end dynamic matching network for multi-view multi-person 3D pose estimation. In \textit{European Conference on Computer Vision}, 477--493. Cham: Springer International Publishing.

\bibitem{Hueser2023} Hueser, T. (2023). JARVIS Annotation Tool. https://github.com/JARVIS-MoCap/JARVIS-AnnotationTool

\bibitem{Iskakov2019} Iskakov, K., Burkov, E., Lempitsky, V., \& Malkov, Y. (2019). Learnable Triangulation of Human Pose. In \textit{Proceedings of the IEEE/CVF International Conference on Computer Vision (ICCV)}, 7718--7727.

\bibitem{Klibaite2025} Klibaite, U., Li, T., Aldarondo, D., Akoad, J. F., Ölveczky, B. P., \& Dunn, T. W. (2025). Mapping the landscape of social behavior. \textit{Cell}, 188(8), 2249--2266.

\bibitem{Karashchuk2021} Karashchuk, P., Rupp, K. L., Dickinson, E. S., Walling-Bell, S., Sanders, E., Azim, E., Brunton, B. W., \& Tuthill, J. C. (2021). Anipose: A toolkit for robust markerless 3D pose estimation. \textit{Cell Reports}, 36(13), 109730.

\bibitem{Lever2006} Lever, C., Burton, S., \& O Keefe, J. (2006).  Rearing on hind legs, environmental novelty, and the hippocampal formation. \textit{Reviews in the Neurosciences}, 17(1/2), 111.

\bibitem{Leonardis2022} Leonardis, E. J., Breston, L., Lucero-Moore, R., Sena, L., Kohli, R., Schuster, L., Barton-Gluzman, L., Quinn,  L. K., Wiles, J., \& Chiba, A. A. (2022). Interactive neurorobotics: Behavioral and neural dynamics of agent interactions. \textit{Frontiers in Psychology}, 13, 897603.

\bibitem{Leonardis2025} Leonardis, E., Nagamori, A., Thanawalla, A., Yang, Y., Park, J., Saunders, H., et al. (2025). Massively Parallel Imitation Learning of Mouse Forelimb Musculoskeletal Reaching Dynamics. \textit{arXiv:2511}.

\bibitem{Maree2024} Maree, L., Afshar, S., Oline, S., Leonardis, E. J., Falkner, A. L., \& Pereira, T. D. (2024). Multi-view triangulation-enabled annotation for multi-animal 3D pose in SLEAP. In \textit{Measuring Behavior}.

\bibitem{Marshall2021} Marshall, J. D., Klibaite, U., Gellis, A., Aldarondo, D. E., Ölveczky, B. P., \& Dunn, T. W. (2021). The PAIR-R24M Dataset for Multi-animal 3D Pose Estimation. \textit{bioRxiv}, 2021.11.23.469743.

\bibitem{Marshall2022} Marshall, J. D., Li, T., Wu, J. H., \& Dunn, T. W. (2022). Leaving flatland: Advances in 3D behavioral measurement. \textit{Current Opinion in Neurobiology}, 73, 102522.

\bibitem{Mathis2018} Mathis, A., Mamidanna, P., Cury, K. M., Abe, T., Murthy, V. N., Mathis, M. W., \& Bethge, M. (2018). DeepLabCut: markerless pose estimation of user-defined body parts with deep learning. \textit{Nature Neuroscience}, 21(9), 1281--1289.

\bibitem{Pereira2020} Pereira, T. D., Shaevitz, J. W., \& Murthy, M. (2020). Quantifying behavior to understand the brain. \textit{Nature Neuroscience}, 23(12), 1537--1549.

\bibitem{Pereira2022} Pereira, T. D., Tabris, N., Matsliah, A., Turner, D. M., Li, J., Ravindranath, S., et al. (2022). SLEAP: A deep learning system for multi-animal pose tracking. \textit{Nature Methods}, 19(4), 486--495.

\bibitem{Waldmann2024} Waldmann, U., Chan, A. H. H., Naik, H., Nagy, M., Couzin, I. D., Deussen, O., ... \& Kano, F. (2024). 3d-muppet: 3d multi-pigeon pose estimation and tracking. \textit{International Journal of Computer Vision}, 132(10), 4235-4252.

\bibitem{Zhang2025} Zhang, C. Y., Yang, Y., Sirbu, A., Abe, E. T., Wärnberg, E., Leonardis, E. J., et al. (2025). MIMIC-MJX: Neuromechanical emulation of animal behavior. \textit{arXiv:2511}.

\end{thebibliography}


\end{document}